% === §2.1 Problemática o necesidad que motiva el proyecto ===
% Redactado por el Investigador — Fase 1

% --- Párrafo 1: Contexto general de la problemática desde el sector de aplicación ---

La educación en ingeniería enfrenta una transformación estructural sin precedentes. La
convergencia entre la cuarta revolución industrial, la masificación de herramientas de
inteligencia artificial generativa y las demandas del sector productivo ha puesto en
evidencia una brecha crítica: los enfoques pedagógicos tradicionales —centrados en la
transmisión expositiva y la evaluación sumativa— resultan insuficientes para desarrollar
las competencias de orden superior que requiere el ingeniero contemporáneo. Investigaciones
recientes en educación superior han documentado que el uso no mediado de asistentes de IA
como ChatGPT o GitHub Copilot por parte de estudiantes de ingeniería puede inhibir el
desarrollo del pensamiento crítico, la argumentación técnica y la capacidad de resolver
problemas auténticos, fenómenos que la literatura educativa agrupa bajo el constructo de
\emph{aprendizaje profundo} \cite{zawacki2019systematic, holmes2022state}. En el
contexto colombiano, esta problemática se agrava por factores estructurales: altas tasas de
deserción en programas de ingeniería, heterogeneidad en los niveles de competencia digital
y matemática de los estudiantes, y una limitada articulación entre los currículos
universitarios y las necesidades del tejido productivo nacional. Los lineamientos del
CONPES 3975 de 2019 y el Plan Nacional de Desarrollo 2022--2026 reconocen explícitamente
la urgencia de modernizar la educación superior mediante la adopción de tecnologías
emergentes, en consonancia con el Objetivo de Desarrollo Sostenible 4 (Educación de
Calidad) y el ODS 9 (Industria, Innovación e Infraestructura). Sin embargo, la mera
incorporación de herramientas de IA en el aula no garantiza una mejora en la calidad del
aprendizaje; se requiere una intervención tecnopedagógica integral que diagnostique,
personalice y evalúe el aprendizaje profundo de manera sistemática.

% --- Párrafo 2: Contexto de los desafíos técnicos desde la perspectiva de la IA + listado de subproblemas ---

Desde la perspectiva de la inteligencia artificial, el panorama actual revela un
desequilibrio significativo. Por un lado, los grandes modelos de lenguaje (LLMs) y las
arquitecturas generativas han alcanzado un nivel de madurez que permite su integración en
entornos educativos con capacidades notables de generación de texto, código y explicaciones
\cite{wang2024survey, park2023generative}. Por otro lado, la mayoría de las aplicaciones
educativas basadas en IA operan como asistentes conversacionales reactivos —``wrappers''
sobre APIs de LLMs— que carecen de un modelo del aprendiz, de memoria pedagógica y de una
orquestación deliberada de estrategias didácticas. Los sistemas tutores inteligentes (ITS)
tradicionales, aunque pedagógicamente fundamentados, se han construido predominantemente
sobre arquitecturas basadas en reglas y árboles de decisión, limitando su adaptabilidad a
la diversidad de trayectorias de aprendizaje en ingeniería. En contraste, el paradigma
emergente de la \emph{IA agéntica} —sistemas compuestos por múltiples agentes autónomos que
perciben, razonan, actúan y colaboran para alcanzar objetivos complejos— ofrece una
oportunidad disruptiva para superar estas limitaciones. Sin embargo, su aplicación
sistemática en educación en ingeniería enfrenta tres subproblemas técnicos interrelacionados
que esta propuesta busca resolver:

\begin{itemize}
    \item \textbf{Subproblema 1 — Diagnóstico y modelado multidimensional del aprendiz:}
    Ausencia de instrumentos y modelos computacionales que permitan diagnosticar de manera
    sistemática, trazable y multidimensional el nivel de aprendizaje profundo —entendido
    como motivación intrínseca, capacidad de argumentación técnica y competencia en
    resolución de problemas auténticos— en estudiantes de ingeniería que interactúan con
    herramientas de inteligencia artificial.

    \item \textbf{Subproblema 2 — Arquitectura agéntica autónoma para personalización del
    aprendizaje:} Carencia de un ecosistema de agentes autónomos basados en IA, validado
    pedagógicamente, que articule de forma coherente funciones de tutoría personalizada,
    retroalimentación adaptativa, generación de problemas auténticos, asistencia en
    programación y depuración, simulación de sistemas ingenieriles y evaluación formativa,
    todo ello orientado al fortalecimiento del aprendizaje profundo en estudiantes de
    ingeniería.

    \item \textbf{Subproblema 3 — Validación empírica y transferencia tecnológica TRL 6/7:}
    Falta de evidencia cuantitativa y cualitativa, obtenida en entornos reales de aula
    universitaria, sobre el impacto causal del ecosistema agéntico en el aprendizaje
    profundo de estudiantes de ingeniería, así como de una ruta de despliegue tecnológico
    que garantice un nivel de madurez TRL 6/7 transferible al entorno educativo y al
    sector productivo colombiano.
\end{itemize}

% --- Párrafo 3: Causas, limitaciones y conceptos del Subproblema 1 ---

El primer subproblema tiene su raíz en una limitación conceptual y metodológica
persistente en la intersección entre IA y educación: la dificultad de operacionalizar el
constructo de \emph{aprendizaje profundo} en métricas computacionales trazables. Si bien
la taxonomía de Marton y Säljö (1976) y el modelo SOLO de Biggs (1999) proporcionan
marcos teóricos sólidos para distinguir entre procesamiento superficial y profundo, su
traducción a señales observables en entornos digitales de aprendizaje sigue siendo un
desafío abierto. Los enfoques actuales de analítica del aprendizaje se concentran en
indicadores superficiales —tasas de clic, tiempo en plataforma, número de intentos— que
no capturan la calidad de la argumentación técnica, la profundidad del razonamiento ni la
transferencia de conocimientos a problemas no estructurados \cite{hooshyar2024systematic}.
En el dominio específico de la ingeniería, esta carencia se magnifica: las tareas
auténticas de ingeniería —diseñar un circuito, depurar un algoritmo, modelar un sistema
físico— requieren formas de razonamiento multimodal (textual, matemático, diagramático)
que los sistemas de diagnóstico actuales no integran. Adicionalmente, la creciente
mediación de LLMs en las actividades académicas introduce un factor de confusión: el
estudiante puede producir respuestas formalmente correctas generadas por IA sin haber
atravesado el proceso cognitivo que conduce al aprendizaje profundo. La ausencia de un
\emph{modelo del aprendiz} computacional que distinga entre desempeño aparente y
comprensión genuina constituye, por tanto, la primera barrera que esta propuesta debe
superar.

% --- Párrafo 4: Causas, limitaciones y conceptos del Subproblema 2 ---

El segundo subproblema se origina en la fragmentación del paisaje tecnológico actual.
Existen asistentes de programación (GitHub Copilot, Cursor), tutores conversacionales
(Khanmigo, Duolingo Max), entornos de simulación (MATLAB, Simulink, LTspice) y
plataformas de evaluación automatizada (Gradescope, CodeRunner). Sin embargo, estas
herramientas operan como islas tecnológicas: ninguna articula en un único ecosistema
agéntico las funciones de diagnóstico del aprendiz, tutoría personalizada, generación de
problemas auténticos, asistencia en programación, simulación de sistemas eléctricos y
electrónicos, y evaluación formativa. La literatura reciente en IA agéntica demuestra que
las arquitecturas multiagente —donde agentes especializados colaboran bajo un protocolo de
orquestación— superan a los sistemas monolíticos en tareas complejas que requieren
descomposición de objetivos, uso de herramientas externas y memoria de largo plazo
\cite{xi2023rise, wang2024survey}. No obstante, el diseño de un ecosistema agéntico para
educación en ingeniería exige resolver desafíos técnicos específicos: (a) la definición de
una ontología de roles agénticos pedagógicamente fundamentada (agente tutor, agente
evaluador, agente generador de problemas, agente simulador); (b) la implementación de un
mecanismo de memoria compartida y contexto pedagógico que preserve la continuidad de la
trayectoria de aprendizaje; (c) la integración de protocolos como el Model Context
Protocol (MCP) para la interacción segura y estandarizada entre agentes, LLMs y
herramientas externas; y (d) la incorporación de principios de \emph{scaffolding}
adaptativo —inspirados en la zona de desarrollo próximo de Vygotsky— que permitan graduar
la asistencia del ecosistema en función del nivel de competencia diagnosticado en el
aprendiz.

% --- Párrafo 5: Causas, limitaciones y conceptos del Subproblema 3 ---

El tercer subproblema refleja una brecha crítica entre la promesa de la IA educativa y su
evidencia empírica. A pesar de la proliferación de estudios sobre IA en educación, las
revisiones sistemáticas más citadas coinciden en señalar que la mayoría de las
intervenciones se evalúan en condiciones controladas de laboratorio, con muestras pequeñas
y sin grupos de control, lo que impide establecer relaciones causales entre el uso de
agentes de IA y mejoras en el aprendizaje profundo \cite{zawacki2019systematic,
hooshyar2024systematic}. En el contexto latinoamericano, esta situación es aún más aguda: son
escasos los estudios que reportan despliegues de sistemas de IA educativa en aulas reales
de ingeniería con seguimiento longitudinal. Paralelamente, la noción de \emph{madurez
tecnológica} (Technology Readiness Level, TRL), ampliamente utilizada en ingeniería para
calificar el grado de preparación de una tecnología para su transferencia al entorno
operativo, ha sido escasamente aplicada al dominio de las tecnologías educativas. La
mayoría de los prototipos de IA para educación en ingeniería se sitúan en TRL 3--4
(prueba de concepto en entorno controlado), sin alcanzar el umbral TRL 6--7 (demostración
en entorno relevante u operativo) que exige la presente convocatoria y que demandan los
potenciales adoptantes —universidades, centros de formación técnica y empresas del sector
productivo—. Cerrar esta brecha requiere un diseño de validación que combine métodos
cuantitativos (diseño cuasiexperimental pre-post con grupo control, analítica del
aprendizaje) y cualitativos (entrevistas semiestructuradas, análisis de la argumentación,
rúbricas de aprendizaje profundo), junto con un plan de despliegue tecnológico que asegure
la escalabilidad, la interoperabilidad y la sostenibilidad del ecosistema más allá del
período de financiación.

% --- Párrafo 6: Conclusión, pregunta de investigación y propuesta de figura (árbol de problemas) ---

Los tres subproblemas descritos no son independientes, sino que conforman una cadena causal
que bloquea el avance de la IA agéntica aplicada a la educación en ingeniería: sin un
diagnóstico multidimensional del aprendiz (SP1), no es posible diseñar trayectorias de
personalización pedagógicamente informadas (SP2); sin una arquitectura agéntica validada,
no puede recogerse evidencia empírica de impacto ni alcanzarse la madurez tecnológica
necesaria para la transferencia (SP3). La presente propuesta aborda esta cadena de manera
integral, articulando fundamentos del aprendizaje profundo, teoría de agentes autónomos,
arquitecturas multiagente basadas en LLMs y metodologías de validación empírica en aula.
En este marco, la pregunta que orienta la investigación es:

\bigskip

\noindent\textbf{¿Cómo desarrollar un ecosistema de agentes autónomos basados en
inteligencia artificial que fortalezca el aprendizaje profundo —motivación intrínseca,
argumentación técnica y resolución de problemas auténticos— en estudiantes de ingeniería,
con un nivel de madurez tecnológica TRL~6/7 transferible al entorno educativo y al sector
productivo colombiano?}

\bigskip

% --- Descripción para la figura de árbol de problemas ---
% NOTA PARA EL DISEÑADOR (agente TikZ): A continuación se describe el contenido conceptual
% para el diagrama de árbol de problemas. El Diseñador deberá implementarlo en TikZ en el
% archivo proposal/sections/diag_arbol_problemas.tex.

\medskip

\noindent\textbf{Propuesta de contenido para la figura del árbol de problemas.}
El diagrama debe representar una estructura jerárquica de tres niveles, siguiendo la
metodología de marco lógico:

\begin{itemize}
    \item \textbf{Raíces (causas estructurales):} (C1) Predominio de pedagogías
    expositivas centradas en la transmisión de contenidos; (C2) Fragmentación del
    ecosistema de herramientas de IA educativa sin orquestación pedagógica; (C3) Escasa
    cultura de validación empírica y transferencia tecnológica en \emph{edtech} para
    ingeniería; (C4) Heterogeneidad en competencias digitales y matemáticas de los
    estudiantes de ingeniería en Colombia.

    \item \textbf{Tronco (problema central):} \emph{``Limitada capacidad del sistema
    educativo en ingeniería para fortalecer el aprendizaje profundo (motivación,
    argumentación técnica y resolución de problemas auténticos) en entornos mediados por
    inteligencia artificial, con tecnologías que no alcanzan madurez suficiente para su
    transferencia al entorno real.''}

    \item \textbf{Ramas (consecuencias):} (E1) Egresados de ingeniería con competencias
    superficiales y dependencia acrítica de herramientas de IA; (E2) Desaprovechamiento
    del potencial de la IA agéntica para personalizar la formación en ingeniería; (E3)
    Baja transferencia de innovaciones en IA educativa desde los laboratorios hacia las
    aulas y el sector productivo; (E4) Desalineación entre los perfiles de egreso y las
    demandas de la Industria 4.0/5.0 en Colombia.
\end{itemize}

El Diseñador debe conectar cada causa (raíz) con el subproblema correspondiente mediante
flechas etiquetadas (C1--C2 $\rightarrow$ SP1; C2--C3 $\rightarrow$ SP2; C3--C4
$\rightarrow$ SP3), y cada consecuencia (rama) con los efectos derivados del problema
central. Se sugiere una paleta cromática institucional (azul UNAL, gris LabIA, verde
GCPDS) y una disposición vertical con las raíces abajo, el tronco al centro y las ramas
arriba.
